% Label style sampler — 8 pages, one label treatment per page
% Fixed pattern: crosshatch 45°+135°, 3.0mm spacing, 0.30pt weight
% Shows squares g1 (DARK), g2 (LIGHT), h1 (LIGHT), h2 (DARK)
% with labels g, h (files) and 1, 2 (ranks)
%
% Label styles:
%  1. All black — uniform black on all squares
%  2. White+black outline (hollow glyph) — contour package
%  3. Black outline box around label — rectangular frame
%  4. Dark badge — white text on dark rounded rect
%  5. White badge (opaque) — black text on white rect, both sq types
%  6. Semi-transparent white badge — white 55% opacity behind black text
%  7. White badge no border — clean white backing, no outline
%  8. Drop shadow — white text with black shadow offset

\documentclass[a4paper]{article}
\usepackage[a4paper, top=25mm, bottom=20mm,
            left=20mm, right=20mm,
            noheadfoot, nomarginpar]{geometry}
\usepackage{tikz}
\usepackage{helvet}
\renewcommand{\familydefault}{\sfdefault}
\usepackage{xcolor}
\usepackage[outline]{contour}
\contourlength{0.7pt}

\pagestyle{empty}

\def\sq{4.9}
\def\lpad{0.30}   % label inset from square edge

% Fixed crosshatch pattern code (used on every page)
\newcommand{\crosshatch}{%
  \foreach \cc in {
    -4.95,-4.53,-4.10,-3.67,-3.25,-2.82,-2.40,-1.97,-1.55,
    -1.12,-0.69,-0.27, 0.15, 0.58, 1.00, 1.43, 1.85, 2.28,
     2.70, 3.13, 3.55, 3.98, 4.40, 4.83}{%
    \draw[black,line width=0.30pt](0,\cc)--(\sq,\cc+\sq);
    \draw[black,line width=0.30pt](0,\cc+\sq)--(\sq,\cc);}%
}

% Draw dark square with crosshatch pattern
\newcommand{\patsquare}[2]{%
  \begin{scope}[shift={(#1,#2)}]
    \fill[white](0,0)rectangle(\sq,\sq);
    \begin{scope}
      \clip(0,0)rectangle(\sq,\sq);
      \crosshatch
    \end{scope}
  \end{scope}
}

% -----------------------------------------------------------------------
% \labelpage{num}{title}{detail}
%   {lblDARKfile}{lblDARKrank}   — label nodes for dark sq file/rank
%   {lblLIGHTfile}{lblLIGHTrank} — label nodes for light sq file/rank
%
% The macro places the 2x2 board then calls the label node commands.
% Each label command receives the anchor position as a pre-set coordinate.
% We use named coordinates: 
%   (Lgfile) = bottom-centre of g1 (dark, bot-left)
%   (Lhfile) = bottom-centre of h1 (light, bot-right)
%   (Lhrank1) = right-centre of h1 (light, bot-right)
%   (Lh2rank) = right-centre of h2 (dark, top-right)
% -----------------------------------------------------------------------
\newcommand{\labelpage}[7]{%
  % #1 num  #2 title  #3 detail
  % #4 dark-file label node  #5 dark-rank label node
  % #6 light-file label node  #7 light-rank label node
  \noindent
  \begin{tikzpicture}[x=1cm, y=-1cm]
    \useasboundingbox(0,0)rectangle(17.1,23.7);

    % Title
    \node[anchor=north west,
          font=\fontsize{16}{20}\selectfont\bfseries,text=black]
      at(0,0){Sample #1 — #2};
    \node[anchor=north west,
          font=\fontsize{10}{13}\selectfont,text=black!55]
      at(0,1.05){#3};

    % Board geometry
    \pgfmathsetmacro{\bx}{(17.1-2*\sq)/2}
    \def\by{2.8}

    % Light squares
    \fill[white](\bx,\by)rectangle(\bx+\sq,\by+\sq);           % g2 top-left
    \fill[white](\bx+\sq,\by+\sq)rectangle(\bx+2*\sq,\by+2*\sq); % h1 bot-right

    % Dark squares with pattern
    \patsquare{\bx+\sq}{\by}          % h2 top-right (DARK)
    \patsquare{\bx}{\by+\sq}          % g1 bot-left  (DARK)

    % Grid lines
    \foreach \k in {0,1,2}{
      \draw[black,line width=0.4pt]
        (\bx+\k*\sq,\by)--(\bx+\k*\sq,\by+2*\sq);
      \draw[black,line width=0.4pt]
        (\bx,\by+\k*\sq)--(\bx+2*\sq,\by+\k*\sq);
    }
    \draw[black,line width=1.4pt]
      (\bx,\by)rectangle(\bx+2*\sq,\by+2*\sq);

    % Named coordinates for label positions
    % g1 (bot-left, DARK): file label bottom-centre
    \coordinate(Lg1file) at (\bx+\sq/2,       \by+2*\sq-\lpad);
    % g2 (top-left, LIGHT): no label
    % h1 (bot-right, LIGHT): file label bottom-centre, rank label right-centre
    \coordinate(Lh1file) at (\bx+3*\sq/2,     \by+2*\sq-\lpad);
    \coordinate(Lh1rank) at (\bx+2*\sq-\lpad, \by+3*\sq/2);
    % h2 (top-right, DARK): rank label right-centre
    \coordinate(Lh2rank) at (\bx+2*\sq-\lpad, \by+\sq/2);

    % Labels — dark squares: #4 (file g on g1), #5 (rank 2 on h2)
    % Labels — light squares: #6 (file h on h1), #7 (rank 1 on h1)
    #4  % file g on g1 (DARK)
    #5  % rank 2 on h2 (DARK)
    #6  % file h on h1 (LIGHT)
    #7  % rank 1 on h1 (LIGHT)

    % Description box
    \draw[black!20,line width=0.5pt,rounded corners=2pt]
      (0,\by+2*\sq+0.55)rectangle(17.1,\by+2*\sq+2.6);
    \fill[black!4,rounded corners=2pt]
      (0,\by+2*\sq+0.55)rectangle(17.1,\by+2*\sq+2.6);
    \node[anchor=north west,font=\fontsize{8}{10}\selectfont\bfseries,text=black!40]
      at(0.3,\by+2*\sq+0.85){LABEL SPECIFICATION};
    \node[anchor=north west,font=\fontsize{10.5}{14}\selectfont,text=black]
      at(0.3,\by+2*\sq+1.3){#3};
    \node[anchor=north west,font=\fontsize{8.5}{11}\selectfont,text=black!50]
      at(0.3,\by+2*\sq+2.0){%
      Dark squares: g1 (bot-left) and h2 (top-right)
      \enspace\textbullet\enspace
      Fixed pattern: crosshatch 45°+135°, 3.0\,mm, 0.30\,pt};

  \end{tikzpicture}
  \newpage
}

% Font shorthand
\newcommand{\lf}{\fontsize{12}{13}\selectfont\bfseries}

\begin{document}

% ============================================================
% Sample 1 — All black labels
% ============================================================
\labelpage{1}{All black labels}{%
  Dark squares: black text, no background \quad
  Light squares: black text, no background}{%
  % g1 file (DARK) — black
  \node[anchor=south,font=\lf,text=black] at(Lg1file){g};}{%
  % h2 rank (DARK) — black rotated
  \node[anchor=west,rotate=180,font=\lf,text=black] at(Lh2rank){2};}{%
  % h1 file (LIGHT) — black
  \node[anchor=south,font=\lf,text=black] at(Lh1file){h};}{%
  % h1 rank (LIGHT) — black rotated
  \node[anchor=west,rotate=180,font=\lf,text=black] at(Lh1rank){1};}

% ============================================================
% Sample 2 — White with black outline (hollow glyph via contour)
% ============================================================
\labelpage{2}{White with black outline (hollow glyph)}{%
  Dark squares: white fill, black perimeter outline (contour)
  \quad Light squares: black text}{%
  \node[anchor=south,font=\lf] at(Lg1file)
    {\contour{black}{\textcolor{white}{g}}};}{%
  \node[anchor=west,rotate=180,font=\lf] at(Lh2rank)
    {\contour{black}{\textcolor{white}{2}}};}{%
  \node[anchor=south,font=\lf,text=black] at(Lh1file){h};}{%
  \node[anchor=west,rotate=180,font=\lf,text=black] at(Lh1rank){1};}

% ============================================================
% Sample 3 — Black outline box (rectangular frame around label)
% ============================================================
\labelpage{3}{Black outline box around label}{%
  Dark squares: black text inside thin black rectangular frame
  \quad Light squares: black text no frame}{%
  \node[anchor=south,font=\lf,text=black,
        draw=black,line width=0.5pt,inner sep=1.2pt] at(Lg1file){g};}{%
  \node[anchor=west,rotate=180,font=\lf,text=black,
        draw=black,line width=0.5pt,inner sep=1.2pt] at(Lh2rank){2};}{%
  \node[anchor=south,font=\lf,text=black] at(Lh1file){h};}{%
  \node[anchor=west,rotate=180,font=\lf,text=black] at(Lh1rank){1};}

% ============================================================
% Sample 4 — Dark badge: white text on dark rounded rect
% ============================================================
\labelpage{4}{Dark badge}{%
  Dark squares: white text on dark grey rounded rect badge
  \quad Light squares: black text on light grey badge}{%
  \node[anchor=south,font=\lf,text=white,
        fill=black!65,rounded corners=1.5pt,inner sep=1.5pt] at(Lg1file){g};}{%
  \node[anchor=west,rotate=180,font=\lf,text=white,
        fill=black!65,rounded corners=1.5pt,inner sep=1.5pt] at(Lh2rank){2};}{%
  \node[anchor=south,font=\lf,text=black,
        fill=black!18,rounded corners=1.5pt,inner sep=1.5pt] at(Lh1file){h};}{%
  \node[anchor=west,rotate=180,font=\lf,text=black,
        fill=black!18,rounded corners=1.5pt,inner sep=1.5pt] at(Lh1rank){1};}

% ============================================================
% Sample 5 — White badge: black text on white opaque rect (both sq types)
% ============================================================
\labelpage{5}{White opaque badge (consistent both square types)}{%
  All squares: black text on white opaque rounded rect — uniform appearance}{%
  \node[anchor=south,font=\lf,text=black,
        fill=white,rounded corners=1.5pt,inner sep=1.5pt] at(Lg1file){g};}{%
  \node[anchor=west,rotate=180,font=\lf,text=black,
        fill=white,rounded corners=1.5pt,inner sep=1.5pt] at(Lh2rank){2};}{%
  \node[anchor=south,font=\lf,text=black,
        fill=white,rounded corners=1.5pt,inner sep=1.5pt] at(Lh1file){h};}{%
  \node[anchor=west,rotate=180,font=\lf,text=black,
        fill=white,rounded corners=1.5pt,inner sep=1.5pt] at(Lh1rank){1};}

% ============================================================
% Sample 6 — Semi-transparent white badge (55% opacity)
% Done by drawing a filled rect first, then placing black text node
% ============================================================
\labelpage{6}{Semi-transparent white badge}{%
  Dark squares: black text, white rect at 55\% opacity behind label
  \quad Light squares: black text no badge}{%
  % g1 file (DARK): manual badge — rect then text
  \fill[white,opacity=0.55]
    (\bx+\sq/2-0.38,\by+2*\sq-\lpad-0.05)
    rectangle(\bx+\sq/2+0.38,\by+2*\sq-\lpad-0.60);
  \node[anchor=south,font=\lf,text=black] at(Lg1file){g};}{%
  % h2 rank (DARK): manual badge rotated area
  \fill[white,opacity=0.55]
    (\bx+2*\sq-\lpad-0.05,\by+\sq/2-0.38)
    rectangle(\bx+2*\sq-\lpad-0.60,\by+\sq/2+0.38);
  \node[anchor=west,rotate=180,font=\lf,text=black] at(Lh2rank){2};}{%
  \node[anchor=south,font=\lf,text=black] at(Lh1file){h};}{%
  \node[anchor=west,rotate=180,font=\lf,text=black] at(Lh1rank){1};}

% ============================================================
% Sample 7 — White badge no border (clean white backing, no outline)
% ============================================================
\labelpage{7}{White badge, no border}{%
  Dark squares: dark grey text on white rect, no border
  \quad Light squares: dark grey text plain}{%
  \node[anchor=south,font=\lf,text=black!80,
        fill=white,inner sep=1.5pt] at(Lg1file){g};}{%
  \node[anchor=west,rotate=180,font=\lf,text=black!80,
        fill=white,inner sep=1.5pt] at(Lh2rank){2};}{%
  \node[anchor=south,font=\lf,text=black!80] at(Lh1file){h};}{%
  \node[anchor=west,rotate=180,font=\lf,text=black!80] at(Lh1rank){1};}

% ============================================================
% Sample 8 — Drop shadow: white text with black offset shadow
% ============================================================
\labelpage{8}{White text with black drop shadow}{%
  Dark squares: white text + black shadow offset 0.5pt
  \quad Light squares: black text no shadow}{%
  % Shadow first (black, offset), then white on top
  \node[anchor=south,font=\lf,text=black]
    at(\bx+\sq/2+0.02,\by+2*\sq-\lpad+0.02){g};
  \node[anchor=south,font=\lf,text=white] at(Lg1file){g};}{%
  \node[anchor=west,rotate=180,font=\lf,text=black]
    at(\bx+2*\sq-\lpad+0.02,\by+\sq/2+0.02){2};
  \node[anchor=west,rotate=180,font=\lf,text=white] at(Lh2rank){2};}{%
  \node[anchor=south,font=\lf,text=black] at(Lh1file){h};}{%
  \node[anchor=west,rotate=180,font=\lf,text=black] at(Lh1rank){1};}

\end{document}
